\scanpage{78}
D'autre part, on a un torseur évident $P_{\unsure{o|0}}$ sur $\handwritten{\underline}{Q}_p$ de groupe $(\handwritten{\hat}{\handwritten{\underline}{Z}})_{\inserted{\handwritten{\underline}{Q}_p}}$, correspondant à l'extension non ramifiée maximale de $\handwritten{\underline}{Q}_p$ ; on\deleted{t} prend l'image \inserted{$P$} de ce torseur: c'est lui.

8.3. Soit $\handwritten{\underline}{H}$ la gerbe noyau de (8.2.1), qui a donc comme lien le groupe diagonalisable $H$ de 8.1. D'après le sorite général (ref) la classe
\margin{\handwritten{$\underline{H}$ est aussi la \uline{gerbe des relèvements}, \uline{pour l'hom.}
\diagram{\begin{tikzcd}[ampersand replacement=\&,column sep=small,row sep=small]
j:\ \underline{G}(\unsure{F|\mathbb{F}}_p) \arrow[r] \arrow[d,phantom,"\simeq"{rotate=90}] \& \underline{\Pi}_p(F_p) \arrow[d,phantom,"\simeq"{rotate=90}] \\
[G] \& [G']
\end{tikzcd}}
déduit de l'hom $\unsure{G|\underline{G}}\longrightarrow G'$, du torseur $P$ ci-dessus.}}
\[
\handwritten{\xi\in}H^{\handwritten{2}}(\handwritten{\underline}{Q}_p,\deleted{\unsure{\underline{G}_m|G_m}}\ H)=\operatorname{Hom}(\handwritten{\underline}{Q},\operatorname{Br}(\handwritten{\underline}{Q}_p))=\operatorname{Hom}(\handwritten{\underline}{Q},\handwritten{\underline}{Q}/\handwritten{\underline}{Z})
\]
de cette gerbe est donnée par
\[
\handwritten{\xi=\unsure{\partial|\delta}(\unsure{\eta|\gamma})}
\]
où
\[
\handwritten{\eta\in}H^{\handwritten{1}}(\handwritten{\underline}{Q}_p,G')
\]
est la classe du torseur $P$. Pour faire le calcul, notons qu'on a un \deleted{sous-groupe}\inserted{quotient} $H'$ de $G_r$, correspondant au sous-groupe de $\deleted{\unsure{Q_p|\underline{Q}_p}}$ $\handwritten{\overline}{\handwritten{\underline}{Q}}_p^{\handwritten{*}}$ \deleted{formé} engendré par les raciens de l'unité(= éléments de torsion) et toutes les raciens de $p$, $p^{1/n}$ (ce groupe\deleted{\unsure{s|x}} est bien invariant par $\operatorname{Gal}(\handwritten{\overline}{\handwritten{\underline}{Q}}_p/\handwritten{\underline}{Q}_p)$ ). On a alors un homomorphisme d'extensions
\margin{\handwritten{\uline{NB} Ce diagramme est induit par un diagramme de \uline{gerbes} (valable sur un corps $k$ parfait quelconque) avec $\underline{H}'=$ gerbe correspondant \unsure{aux|à} la sous-catégorie de $\operatorname{Fcriso}(F_p)$ formé\unsure{e|s} des objets qui deviennent isomorphes sur une extension \uline{finie} $k'$ de $k=F_p$ à une somme d'objets $E_{r,s}$ ; et $H$ s'identifie à la \uline{gerbe des relèvements} de $P_o$ pour l'hom. $j:\underline{H}'\longrightarrow[\hat{\underline{Z}}]$ associé à $H'\longrightarrow(\hat{\underline{Z}})_{\mathbb{Q}_p}$ ...}}
\diagram{\begin{tikzcd}[ampersand replacement=\&]
0 \handwritten{\arrow[r]} \& H \handwritten{\arrow[r]} \handwritten{\arrow[d]} \& G \handwritten{\arrow[r]} \handwritten{\arrow[d]} \& G' \handwritten{\arrow[r]} \handwritten{\arrow[d]} \& 0 \\
0 \handwritten{\arrow[r]} \& H \handwritten{\arrow[r]} \& H' \handwritten{\arrow[r]} \& \handwritten{(\hat{\underline{Z}})_{\mathbb{Q}_p}} \handwritten{\arrow[r]} \& 0\ ,
\end{tikzcd}}
qui montre qu'on \deleted{\unsure{suffit de calculer|suffit de calculer xxx}} a
\[
\handwritten{\xi=\partial(\eta_o)},
\]
où
\[
\handwritten{\eta_o\in}H^{\handwritten{1}}(\handwritten{\underline}{Q}_p,\handwritten{\hat}{\handwritten{\underline}{Z}})
\]
est l'image de $\handwritten{\eta}$ par $G'\handwritten{\longrightarrow}\handwritten{\hat}{\handwritten{\underline}{Z}}$. On constate que le composé de (8.2.2) avec l'homomorphisme précédent est l'identité de $\handwritten{\hat}{\handwritten{\underline}{Z}}$, donc $\handwritten{\eta_o}$ est la classe du torseur canonique $P_o$ de la fin de 8.2. A l'aide\deleted{\unsure{x|s}} de ceci, il faut vérifier que la classe $\handwritten{\xi}$ correspond (par l'isomorphisme
\[
\operatorname{Br}(\handwritten{\underline}{Q}_p)\handwritten{\simeq}\handwritten{\underline}{Q}_{\deleted{\handwritten{\unsure{p|\ell}}}}/\handwritten{\underline}{Z}_{\deleted{\handwritten{\unsure{p|\ell}}}}
\]
du corps de classe local, à l'homomorphisme canonique $\handwritten{\mathbb{Q}\longrightarrow\mathbb{Q}/\mathbb{Z}}$
% archive sheet number: 77
